% RESEARCH_DRAFT_SHA256: 86e8efacb4829e31c3b01b2c6d532a93c4b09167dca67ac88fb07352625f0c8d
% Layout and source-reference wrappers expanded for the web renderer.
\documentclass[11pt]{article}
\usepackage[margin=25mm]{geometry}
\usepackage{amsthm}
\usepackage{fontspec,unicode-math}
\setmainfont{Latin Modern Roman}
\setmathfont{Latin Modern Math}
\usepackage{amsmath,amssymb,mathtools,xurl,hyperref}
\hypersetup{colorlinks=true,urlcolor=blue,linkcolor=blue}
\setcounter{secnumdepth}{0}
\setlength{\parindent}{0pt}
\setlength{\parskip}{5pt}
\setlength{\emergencystretch}{3em}
\newcommand{\sref}[2]{\hyperlink{source:#1:#2}{[S#2]}}
\newcommand{\cataloguescope}[1]{\paragraph{Source record.}#1}
\newcommand{\floor}[1]{\left\lfloor#1\right\rfloor}
\newtheorem{proposition}{Proposition}
\newtheorem{lemma}[proposition]{Lemma}
\begin{document}
\paragraph{Research manuscript.} OpenAI \texttt{gpt-6-astra}, reasoning \texttt{ultra}; 11 October 2026.
\paragraph{Preserved source body.} The complete catalogue mathematical text and sources follow with layout and source links expanded, including source qualifications and its historical preparation notes. New research begins at the section headed Research attempt.

% UnsolvedMath ID 30002239; source catalogue code OWR-12177-008
\section{Polynomial Freiman–Ruzsa Conjecture over Binary Vector Spaces}\label{30002239}
{\small UnsolvedMath ID 30002239 · Computer Science}
\subsection{Definitions and mathematical statement}

\paragraph{Shared notation.}Standard number systems are $\mathbb N,\mathbb Z,\mathbb Q,\mathbb R,\mathbb C$. Specialized notation and index ranges follow the cited source and the definitions below.


In $\mathbb F_2^n$, $A+A=\{a+b:a,b\in A\}$ uses coordinatewise addition modulo two. The doubling constant is $|A+A|/|A|$; $\operatorname{span}(A)$ is the smallest linear subspace containing $A$. Polynomial losses in $K$ mean fixed powers independent of $n$. The ordinary small-doubling hypothesis assumes $A\ne\varnothing$ and $K\ge1$.
\paragraph{Record scope and qualifications.}
The catalogue title is retained as metadata. Named certificates, contracts, bounds or frameworks in that title are not asserted as results of the original source extract.
The supplied snapshot labels this record open, but its literature assessment cites the solution of the characteristic-two Marton/Polynomial Freiman-Ruzsa conjecture for the intended nonempty-set formulation. It is retained as a historical source formulation, not asserted to remain an open conjecture.

Source-record limitation: The supplied assessment flags a survey extraction with finite-field solutions; importing its binary PFR formulation as an open target would conflict with that supplied status evidence. The supplied wording does not exclude $A=\varnothing$: then $A^{\prime}=\varnothing$ and $|\operatorname{span}(A^{\prime})|=1$, so its displayed span bound cannot hold. The solved-literature comparison applies to the ordinary nonempty-set formulation; the source statement is preserved.
\paragraph{Statement recorded in the supplied source.}
If $A\subseteq\mathbb F_2^n$ satisfies $|A+A|\leq K|A|$, is there a subset $A^{\prime}\subseteq A$ with $|A^{\prime}|\geq |A|/\operatorname{poly}(K)$ and $|\operatorname{span}(A^{\prime})|\leq\operatorname{poly}(K)|A^{\prime}|$? \href{https://doi.org/10.4171/owr/2012/54}{[S2]}
\paragraph{Source status.}
The supplied catalogue labels this source record open; that metadata was checked against the supplied snapshot on 2026-10-11. The label is preserved for provenance and is not an independent verification of present open status or mathematical correctness. The catalogue literature assessment is dated 2026-08-21.
\subsection{Short English statement}
Relate a small-doubling binary set to a large subset with comparably small linear span.
\subsection{Sources}
\begin{itemize}
\item[S1] Ulam AI and the UnsolvedMath Contributors, UnsolvedMath: A Curated Collection of Open Mathematics Problems, record 30002239 (OWR-12177-008). Catalogue curation and metadata: CC BY 4.0. The source extract is retained; definitions, English summaries and qualifications are editorial additions.
\url{https://huggingface.co/datasets/ulamai/UnsolvedMath}.
\item[S2] Complexity Theory (2013), Oberwolfach Reports, DOI 10.4171/owr/2012/54, pp. 3267--3304
\url{https://doi.org/10.4171/owr/2012/54}.
\item[S3] Gowers--Green--Manners--Tao, Marton conjecture in characteristic two
\url{https://arxiv.org/abs/2311.05762}.
\item[S4] Kazemi Moghadam, odd-characteristic finite-field case
\url{https://arxiv.org/abs/2512.04433}.
\item Source-specific Oberwolfach Reports licences, including CC BY-SA 4.0 where applicable, are retained; see THIRD\_PARTY\_NOTICES.md.
\end{itemize}
\paragraph{Preparation and review.}
The original source and UnsolvedMath contributors retain attribution for the source extract. Codex (GPT-6; exact serving revision not exposed) assisted with definitions, English summaries, scope notes and formatting. No new human mathematical authorship or independently verified solution is asserted. Permanent source ID 30002239 and code OWR-12177-008 checked against the complete supplied registry and current repository. Distinct upstream records, including repeated questions and broad or damaged extracts, are retained for complete coverage of the source CS category. A repeated record is not a new independent conjecture. Complete mathematical-equivalence and current-literature review has not been performed. Preparation and final review follow contribution prompts C and D, with source gaps disclosed. Structural and rendering checks do not establish a new theorem or open status.
\label{end:30002239}

\section{Research attempt}
\subsection{Literature, source context and scope}
The full original catalogue body and source qualifications are preserved above. The relevant primary workshop passages were read in the cached publisher report. Historical claims and subsequent attributed results are distinguished from the mathematical work below. A restricted example, standard lemma or known theorem is not counted as new progress on the entire original target. Missing assumptions and repeated or rhetorical extracts are not silently turned into different solved problems.

\subsection{An attributed resolution and a direct endpoint calculation}
The preserved record already flags the subsequently proved characteristic-two polynomial Freiman--Ruzsa theorem and the empty-set issue. The 2023 primary paper proves Marton's conjecture in the nonempty finite setting. Its covering formulation supplies an absolute constant $C$ such that a nonempty $A\subseteq\mathbb F_2^n$ with $|A+A|\le K|A|$ is covered by at most $N=2K^C$ cosets of a subgroup $H$ with $|H|\le|A|$. We use this as an attributed external theorem, not as a proof produced or fully verified in this attempt.
\begin{proposition}
The stated covering conclusion implies that some nonempty $A'\subseteq A$ has $|A'|\ge|A|/N$ and $|\operatorname{span}(A')|\le2N|A'|$. Independently of that theorem, when $K=1$, $A$ itself is an affine subspace and $|\operatorname{span}(A)|\le2|A|$.
\end{proposition}
\begin{proof}
At least one covering coset $t+H$ meets $A$ in at least $|A|/N$ elements; take this intersection for $A'$. Its linear span is contained in $H+\operatorname{span}\{t\}$, a subgroup of size at most $2|H|$. The hypotheses give $2|H|\le2|A|\le2N|A'|$.
For the endpoint, choose $a_0\in A$ and put $B=A+a_0$. Then $0\in B$, so $B\subseteq B+B$. In characteristic two, $|B+B|=|A+A|\le|A|=|B|$, hence $B+B=B$. Closure under addition, the presence of zero, and the identity $b+b=0$ show that $B$ is a subgroup, equivalently a linear subspace. Thus $A=a_0+B$, and its span is contained in $B+\operatorname{span}\{a_0\}$ of size at most $2|B|=2|A|$.
\end{proof}
The attempted elementary extension from $|A+A|=|A|$ to general small doubling stops at a precise step: $B\subseteq B+B$ no longer forces equality once $K>1$. For example $B=\{00,10,01\}\subseteq\mathbb F_2^2$ has $|B+B|=4$, so it is not closed under addition despite doubling $4/3$. The elementary endpoint proof therefore supplies no polynomial general bound.
\subsection{Verification and remaining gaps}
All nonempty subsets of $\mathbb F_2^3$ were enumerated to check that the doubling-one sets are affine subspaces and satisfy the span bound. The three-element obstruction was checked exactly. The general theorem is credited to its authors; only its short consequence and the endpoint calculation are proved here. No new solution or independent full proof of polynomial Freiman--Ruzsa is claimed. The original catalogue and source status remain unchanged, and this attempt records zero new progress.

\section{References and source attribution}
Primary workshop report: \url{https://doi.org/10.4171/owr/2012/54}. Relevant passages read: 2012/54 pp. 3279--3280, the historical binary polynomial Freiman--Ruzsa discussion.

\url{https://arxiv.org/abs/2311.05762}. Read primary abstract: characteristic-two Marton conjecture proved; result attributed to Gowers, Green, Manners and Tao.

\url{https://arxiv.org/html/2311.05762v2}. Read primary introduction and covering formulation with polynomially many cosets and subgroup size at most |A|; full theorem proof not independently verified.

All retrievals were checked on 11 October 2026. Original question attribution remains in the preserved Sources section. External theorems are attributed and not claimed as new discoveries or fully verified merely by citation.


\paragraph{Integration review.} Codex (GPT-6; exact serving revision not exposed) reviewed the written arguments and expanded layout and source-reference wrappers. Repository metadata and the production web renderer were checked using Python 3.13 and JavaScript with MathJax 3.2.2. This model-assisted review is not independent human proof verification; no PDF compilation or formal proof checking is asserted.

\section{Completion Estimate}
\noindent\textbf{COMPLETION: 0\%}
No substantiated new progress toward the entire original target is claimed. The worked results and failed route are delimited above; a known result, illustrative model, or auxiliary obstruction is not counted as original-target completion.

\section{Final status and attribution}
\textbf{Attempt status: unresolved.} This describes this research attempt, not a new ruling on the source catalogue's status or on external literature claims.
The mathematics and this self-audit were generated by OpenAI \texttt{gpt-6-astra}, reasoning \texttt{ultra}, on 11 October 2026. Codex (GPT-6; exact serving revision not exposed) orchestrated the batch, reviewed the written mathematical arguments, and contributed formatting/repository review. No human mathematical authorship, independent human proof checking, or formal verification is claimed. The historical preparation paragraphs in the reproduced catalogue body concern that body only.

Contribution prompt A was applied to the complete definition and supplied source record, followed by the adversarial checks of prompt B. Local reproduction files are \texttt{batch019/checks.py}, \texttt{computations.json}, and the per-ID audit JSON. They were not public links at preparation time. Numerical checks are finite and do not prove unrestricted results. No LaTeX engine was found locally; compilation was not performed by this research agent.

\end{document}
